% Chapter "Computing RH: the barrier method in Thorin form". Intended place: Volume II, Part I, after the chapter on the de Bruijn Newman flow
% (vol/v2_ch04.tex, Section ch:dbn) and before the Thorin phase diagram (Chapter chap:phase). Numbers: code/nc_barrier_check.py, data/nc_barrier_check.json.

This chapter states the barrier method of Polymath15 \cite{Polymath}, as used by Platt and Trudgian \cite{PT21}, translates it into the language of Thorin measures, and runs a small barrier check of our own. The method proves upper bounds for the de Bruijn Newman constant $\Lambda$ by combining three things: a verification of RH to a height, an analytic zero free region for $H_{t_0}$ at large heights, and a barrier, a window of heights through which no zero of $H_t$ can pass at any time $t\in[0,t_0]$. With it Polymath15 proved $\Lambda\le0.22$ and Platt and Trudgian $\Lambda\le0.2$ (Section~\ref{ch:dbn} and Table~\ref{tab:ch21:hist}).

The main results are these. The barrier criterion is stated precisely as Polymath15 state it, with the role of each hypothesis (Theorem~\ref{barrier:thm:pm}, cited). In Thorin language (Section~\ref{barrier:sec:thorin}) its conclusion is that the shifted clock $\varphi_{t_0,\alpha_0}$ with $\alpha_0=\frac12+\frac{y_0}2$ is a GGC transform, so that by Theorem~\ref{rz:thm:dbn} the reciprocal $R_t$ of $H_t$ is a GGC transform for every $t\ge t_0+\frac12y_0^2$ (Proposition~\ref{barrier:prop:ggc}, proved); the barrier becomes the statement that no atom of the Thorin measure of $R_t$, in the chart of squares, crosses a curvilinear window bounded by four confocal parabolas with focus at the origin (Lemma~\ref{barrier:lem:quad} and Proposition~\ref{barrier:prop:window}, proved); and a positive density $v_t=\re(\xi_t'/\xi_t)$ on the barrier rectangle at every time is a sufficient certificate (Proposition~\ref{barrier:prop:density}, proved). Section~\ref{barrier:sec:check} runs the barrier at the modest height $X/2=497.4$ for $t\in[0,0.2]$ and $y_0=0.2$ in floating point, and certifies Proposition~\ref{barrier:prop:check} in ball arithmetic for every $t\in[0,0.2]$. Section~\ref{barrier:sec:limits} explains why the method, in either language, cannot reach $t=0$.

For the path to RH of Section~\ref{ov1:sec:path} the barrier method is the source we know of an analytic tail below the basepoint one, and the overture names its limit: what separates $\Lambda\le0.2$ from RH is the existence of an analytic region beyond the barrier at positive time and its absence at time zero. The Thorin form makes that statement exact but does not move it. Every point the method certifies lies on a level set $t+2(\alpha-\frac12)^2=c$ with $c>0$, and the corner $(0,\frac12)$ of the phase diagram (Theorem~\ref{cp:thm:phase}) is reached only in the limit, where the analytic hypothesis becomes a zero free region for $\zeta$ itself.

\section{The barrier argument}\label{barrier:sec:argument}

We use the normalisation of Section~\ref{sec:ch21:flow}: $H_t(z)=\int_0^\infty e^{tu^2}\Phi_{\rm dBN}(u)\cos(zu)\,du$, $\xi_t(s)=8H_t(-2i(s-\frac12))$, and a zero $x+iy$ of $H_t$ is the zero $\frac12-\frac y2+\frac{ix}2$ of $\xi_t$. The zeros of $H_t$ are symmetric under $z\mapsto-z$ and $z\mapsto\bar z$, so it suffices to consider $x>0$, $y\ge0$. By de Bruijn's strip theorem (Theorem~\ref{tf:thm:heat}(c) with $r=1$) every zero of $H_t$ satisfies $y^2\le\max(1-2t,0)$ for $t\ge0$.

\subsection{De Bruijn's theorem and the canopy}

The method rests on one consequence of the strip theorem. If at a time $t_0\ge0$ no zero of $H_{t_0}$ satisfies $y>y_0$, then all zeros of $H_{t_0}$ lie in $|\im z|\le y_0$, and by the strip theorem started at $t_0$ the zeros of $H_t$ satisfy $y^2\le\max(y_0^2-2(t-t_0),0)$, so they are real for $t\ge t_0+\frac12y_0^2$. Hence
\begin{equation}\label{barrier:eq:dB}
\Lambda\le t_0+\tfrac12y_0^2\quad\text{whenever $H_{t_0}$ has no zero with } y>y_0 .
\end{equation}
This is \cite[Theorem 13]{deBruijn} in the form of \cite{Polymath}, and in the language of the phase diagram it is Theorem~\ref{cp:thm:phase}(ii) at the point $(t_0,\frac12+\frac{y_0}2)$. Since $y\le\sqrt{1-2t_0}$ in any case, the hypothesis of \eqref{barrier:eq:dB} is that the \emph{canopy}
\[
\mathcal C(t_0,y_0)=\{x+iy:\ x\ge0,\ y_0<y\le\sqrt{1-2t_0}\}
\]
contains no zero of $H_{t_0}$. The canopy is unbounded in $x$, and no finite computation can clear it. The barrier method splits it at a height $X$ into a part to the left, cleared by the verification of RH and by the dynamics of the zeros, and a part to the right, cleared by an analytic estimate.

\subsection{The criterion of Polymath15}

We state the criterion as it is printed in \cite{Polymath} (Theorem 1.2 there, read in the arXiv version and its HTML rendering), with $\zeta$ and $H_t$ in our notation.

\begin{theorem}[Polymath15, upper bound criterion \cite{Polymath}]\label{barrier:thm:pm}
Let $0<t_0\le\frac12$, $X>0$ and $0<y_0\le1$, and suppose:
\begin{enumerate}
\item[(i)] (numerical verification of RH at time $0$) there is no zero $\zeta(\sigma+iT)=0$ with $\frac{1+y_0}2\le\sigma\le1$ and $0\le T\le X/2$;
\item[(ii)] (asymptotic zero free region at time $t_0$) there is no zero $H_{t_0}(x+iy)=0$ with $x\ge X+\sqrt{1-y_0^2}$ and $y_0\le y\le\sqrt{1-2t_0}$;
\item[(iii)] (barrier at intermediate times) there is no zero $H_t(x+iy)=0$ with $X\le x\le X+\sqrt{1-y_0^2}$, $\sqrt{y_0^2+2(t_0-t)}\le y\le\sqrt{1-2t}$ and $0\le t\le t_0$.
\end{enumerate}
Then $\Lambda\le t_0+\frac12y_0^2$.
\end{theorem}

Polymath15 describe the roles of the hypotheses as follows: (i) prevents zeros from reaching the canopy from an initial position to the left of the barrier, (ii) prevents zeros from lying in the canopy to the right of the barrier, and (iii) prevents zeros from starting to the right of the barrier and reaching the canopy to the left of it. The proof combines \eqref{barrier:eq:dB} with the continuity of the zeros in $t$ and the strip theorem applied on subintervals; we do not reproduce it. In practice \cite{Polymath} replace the region of (iii) by the larger box
\begin{equation}\label{barrier:eq:box}
X\le x\le X+1,\qquad y_0\le y\le1,\qquad 0\le t\le t_0,
\end{equation}
and they take $t_0=y_0=0.2$ and $X=6\times10^{10}+83952-0.5$, a value near the limit of the verifications of RH available to them, shifted so that a partial Euler product over the primes up to $11$ is large on the barrier and $|H_t|$ stays away from zero there. This gives $\Lambda\le0.22$ (Theorem 1.1 of \cite{Polymath}). By their account (ii) is the expensive hypothesis, verified with an effective Riemann Siegel type approximation of $H_t$ in the region $0<t\le\frac12$, $0\le y\le1$, $x\ge200$, and (iii) is verified numerically and quickly by comparison. Platt and Trudgian take the second row of the table of conditional bounds of \cite{Polymath}, which needs RH to height $2.51\times10^{12}$, and their verification to $3\times10^{12}$ gives $\Lambda\le0.2$ \cite{PT21}; Section~\ref{cp:sec:cost} records the cost of further rows. The barrier idea itself was proposed in the seventh thread of the project, as a region $0\le t\le t_0$, $X\le x\le X+1$, $y\ge y_0$ that prevents zeros from flowing into the canopy at time $t_0$, assuming they were not already there at time $0$ \cite{NCPolymath7}.

\begin{remark}\label{barrier:rem:dynamics}
The mechanism behind (iii) is the dynamics of Section~\ref{sec:ch21:flow}. Forward in time the zeros repel, $\dot z_k=2\sum_{j\ne k}(z_k-z_j)^{-1}$, a complex zero is pulled towards the real axis by the strip theorem, and zeros are continuous in $t$ (Hurwitz's theorem; at a collision Lemma~\ref{ft:lem:split} describes the splitting). Heuristically, a zero that lies in the canopy at time $t_0$ and to the left of the barrier was higher at earlier times, as the bound on the highest zero suggests: the height $y(t)$ of the highest zero satisfies $y(t)^2\ge y(t_0)^2+2(t_0-t)$ by the strip theorem applied on $[t,t_0]$; so it was either to the left of the barrier at time $0$, which (i) excludes, or it crossed the window $X\le x\le X+1$ at some intermediate time at a height above $\sqrt{y_0^2+2(t_0-t)}$, which (iii) excludes. The width $\sqrt{1-y_0^2}\le1$ of the window is what prevents a zero from jumping over it: by the strip theorem a zero cannot move far horizontally without passing through it, a point that \cite{Polymath} make precise in their proof.
\end{remark}

\subsection{The rectangle in the variable of \texorpdfstring{$\xi_t$}{xi t}}

The box \eqref{barrier:eq:box} becomes a rectangle in $s$. A zero $x+iy$ of $H_t$ with $y\ge0$ gives the zero $\frac12-\frac y2+\frac{ix}2$ of $\xi_t$, and by $\xi_t(1-s)=\xi_t(s)$ (Lemma~\ref{cp:lem:structure}) and conjugation also $\frac12+\frac y2+\frac{ix}2$. So the box at time $t$ is free of zeros exactly when $\xi_t$ has no zero in
\begin{equation}\label{barrier:eq:B}
\mathcal B(X,y_0)=\big[\tfrac12+\tfrac{y_0}2,\,1\big]\times\big[\tfrac X2,\,\tfrac X2+\tfrac12\big].
\end{equation}
The right edge $\re s=1$ carries no zero for $t\ge0$: at $t=0$ because $\zeta(1+i\theta)\ne0$, and for $t>0$ because every zero has $|\beta-\frac12|\le\frac12\sqrt{1-2t}<\frac12$ (Lemma~\ref{cp:lem:right}). The barrier at a given time is therefore the vanishing of the winding number of $\xi_t$ around $\partial\mathcal B$, together with the absence of zeros on $\partial\mathcal B$; and since zeros move continuously, a winding number of $0$ at the times of a grid together with a Rouch\'e bound between grid times (hypothesis (B2) of Section~\ref{sec:ft:persist}) gives (iii) for all $t\in[0,t_0]$.

\section{The barrier in Thorin form}\label{barrier:sec:thorin}

\subsection{The conclusion as a GGC statement}

\begin{proposition}\label{barrier:prop:ggc}
Under the hypotheses of Theorem~\ref{barrier:thm:pm}, let $\alpha_0=\frac12+\frac{y_0}2$. Then:
\begin{enumerate}
\item[(a)] the shifted clock $\varphi_{t_0,\alpha_0}(u)=\xi_{t_0}(\alpha_0)/\xi_{t_0}(\alpha_0+\sqrt u)$ is the Laplace transform of a GGC, whose Thorin measure in the chart $\theta=\sqrt z$ is the Poisson sum of Theorem~\ref{cp:thm:transfer};
\item[(b)] $(t,\alpha)\in\mathcal G$ for every $t\ge t_0$ and $\alpha\ge\frac12$ with $(\alpha-\frac12)^2+\frac t2\ge\frac14y_0^2+\frac{t_0}2$;
\item[(c)] for every $t\ge t_0+\frac12y_0^2$ the reciprocal $R_t(u)=H_t(0)/H_t(2i\sqrt u)$ is the Laplace transform of a GGC with Thorin measure $\sum_k\delta_{(z_k(t)/2)^2}$, $z_k(t)>0$ the zeros of $H_t$.
\end{enumerate}
\end{proposition}

\begin{proof}
Theorem~\ref{barrier:thm:pm} is proved in \cite[Proposition 3.3]{Polymath} through \eqref{barrier:eq:dB}, that is, by showing that $H_{t_0}$ has no zero with $y>y_0$; in the variable of $\xi_{t_0}$ this says that every zero has $\re\rho\le\frac12+\frac{y_0}2=\alpha_0$. Theorem~\ref{cp:thm:transfer} then gives (a). Part (b) is Theorem~\ref{cp:thm:phase}(i) from the point $(t_0,\alpha_0)$, and (c) is the case $\alpha=\frac12$ of (b) combined with Theorem~\ref{rz:thm:dbn}, since $\Lambda\le t_0+\frac12y_0^2$.
\end{proof}

With $t_0=y_0=0.2$ this is Corollary~\ref{cp:cor:known}(b). The proposition only renames the conclusion; what this section adds is the reading of the hypotheses.

\subsection{Atoms in the chart of squares}

Let $t\ge0$. By Lemma~\ref{ki:lem:dict}, applied to $F=H_t$ (whose kernel $e^{tu^2}\Phi_{\rm dBN}$ satisfies part (b) of the lemma), the function $R_t$ is an infinite product of factors $(1+4u/w_k^2)^{-1}$ over one representative $w_k$ of each pair $\pm w_k$ of zeros of $H_t$, and we choose $w_k$ with $\re w_k>0$, or $\im w_k>0$ if $\re w_k=0$ (which does not occur, since $H_t>0$ on the imaginary axis). We call
\[
\mu_t=\sum_k\delta_{w_k(t)^2/4}
\]
the \emph{formal Thorin measure} of $R_t$. It is a positive measure on $(0,\infty)$, and then the Thorin measure of $R_t$, exactly when all zeros are real, that is when $t\ge\Lambda$; a nonreal zero gives a pair of atoms $w^2/4$ and $\bar w^2/4$ off the half line, and then $R_t$ is not a GGC transform (Lemma~\ref{ki:lem:dict}(a)). In the language of $\mu_t$, a collision of two real zeros as $t$ decreases is a collision of two atoms on $(0,\infty)$ after which a conjugate pair of atoms leaves the half line, and RH is the statement that no atom of $\mu_t$ leaves $(0,\infty)$ for any $t>0$.

The map $w\mapsto w^2$ carries the rectangular coordinates of $w=x+iy$ to the parabolic coordinates of $w^2$. Lemma~\ref{ki:lem:parab} treats the horizontal lines; the vertical lines behave in the same way.

\begin{lemma}[The barrier as a parabolic quadrilateral]\label{barrier:lem:quad}
For $b>0$ let $\mathcal P_b=\{\xi+i\eta:\ \xi=\eta^2/(4b^2)-b^2\}$, the right opening parabola of Lemma~\ref{ki:lem:parab}, and for $a>0$ let $\mathcal Q_a=\{\xi+i\eta:\ \xi=a^2-\eta^2/(4a^2)\}$, a left opening parabola. All of them have their focus at the origin, and $\mathcal P_b$ meets $\mathcal Q_a$ at right angles at the points $(a\pm ib)^2$. The map $w\mapsto w^2$ is a bijection from the box $\{x+iy:\ a_1\le x\le a_2,\ b_1\le y\le b_2\}$, $0<a_1<a_2$, $0<b_1<b_2$, onto the closed curvilinear quadrilateral in the upper half plane bounded by arcs of $\mathcal P_{b_1}$, $\mathcal P_{b_2}$, $\mathcal Q_{a_1}$ and $\mathcal Q_{a_2}$. In particular the box \eqref{barrier:eq:box} goes to the quadrilateral $\mathcal W(X,y_0)$ bounded by $\mathcal P_{y_0}$, $\mathcal P_1$, $\mathcal Q_X$ and $\mathcal Q_{X+1}$, and the canopy $\mathcal C(t_0,y_0)$ goes to the region of the closed upper half plane strictly to the left of $\mathcal P_{y_0}$ and on or to the right of $\mathcal P_{\sqrt{1-2t_0}}$.
\end{lemma}

\begin{proof}
For $w=a+iy$ with $a>0$ fixed, $w^2=(a^2-y^2)+2iay$, so $\eta=2ay$ and $\xi=a^2-\eta^2/(4a^2)$: the vertical line goes to $\mathcal Q_a$. The horizontal line $w=x+ib$ goes to $\mathcal P_b$ by the proof of Lemma~\ref{ki:lem:parab}. A parabola $\xi=\eta^2/(4p)-p$ has its focus at $(0,0)$ for either sign of $p$, which gives the common focus. The map is conformal on $\re w>0$, so the images of the orthogonal lines $x=a$ and $y=b$ meet at right angles at $(a+ib)^2$, and at $(a-ib)^2$ by symmetry. The map is injective on the open right half plane, where it is a conformal bijection onto $\cut$, so the box goes onto the region bounded by the images of its four sides, which lies in the upper half plane since $\eta=2xy>0$ there. The statement for the canopy follows in the same way from the horizontal lines alone, together with $x\ge0$, which is $\eta\ge0$.
\end{proof}

The atoms of $\mu_t$ live in the scaled chart $w^2/4$, and scaling by $\frac14$ maps confocal parabolas to confocal parabolas; we keep the chart $w^2$ for the geometry, as Lemma~\ref{ki:lem:parab} does. In this chart de Bruijn's theorem reads: if at time $t_0$ every atom lies on or to the right of $\mathcal P_{y_0}$ (or on the half line), then at time $t$ every atom lies on or to the right of $\mathcal P_{\sqrt{y_0^2-2(t-t_0)}}$, and the parabolas collapse onto $[0,\infty)$ at $t=t_0+\frac12y_0^2$ (Corollary~\ref{ki:cor:parab} is the case $t_0=0$, $y_0=1$).

\begin{proposition}[The barrier as a window]\label{barrier:prop:window}
Write $\mathcal W_t$ for the image under $w\mapsto w^2$ of the region of hypothesis (iii) of Theorem~\ref{barrier:thm:pm} at time $t$, the quadrilateral bounded by $\mathcal Q_X$, $\mathcal Q_{X+\sqrt{1-y_0^2}}$, $\mathcal P_{\sqrt{y_0^2+2(t_0-t)}}$ and $\mathcal P_{\sqrt{1-2t}}$. Then the hypotheses of Theorem~\ref{barrier:thm:pm} are equivalent to the following statements about the formal Thorin measures:
\begin{enumerate}
\item[(i$'$)] $\mu_0$ has no atom in the image of $\{0\le x\le X,\ y_0\le y\le1\}$, the part of the closed region between $\mathcal P_{y_0}$ and $\mathcal P_1$ to the left of $\mathcal Q_X$, scaled by $\frac14$;
\item[(ii$'$)] $\mu_{t_0}$ has no atom in the part of the closed region between $\mathcal P_{y_0}$ and $\mathcal P_{\sqrt{1-2t_0}}$ on or to the right of $\mathcal Q_{X+\sqrt{1-y_0^2}}$, scaled by $\frac14$;
\item[(iii$'$)] for every $t\in[0,t_0]$, $\mu_t$ has no atom in $\frac14\mathcal W_t$.
\end{enumerate}
Under (i$'$) to (iii$'$), no atom of $\mu_{t_0}$ lies strictly to the left of $\mathcal P_{y_0}$ off the half line, and $\mu_t$ is a positive measure on $(0,\infty)$ for $t\ge t_0+\frac12y_0^2$.
\end{proposition}

\begin{proof}
The atoms of $\mu_t$ are the points $w^2/4$ with $w$ a zero of $H_t$ in the open right half plane; since $\bar w$ is then also such a zero, the set of atoms is invariant under conjugation, and it suffices to consider the atoms with $\eta\ge0$, which come from the zeros with $x>0$, $y\ge0$. Lemma~\ref{barrier:lem:quad} maps each region of Theorem~\ref{barrier:thm:pm} bijectively onto the region named. For (i$'$), hypothesis (i) is the absence of zeros $\frac12+\frac y2+\frac{ix}2$ of $\xi$ with $y_0\le y\le1$ and $0\le x\le X$, which is the absence of zeros of $H_0$ in that box. The last sentence is Theorem~\ref{barrier:thm:pm} with \eqref{barrier:eq:dB} and the definition of $\mu_t$.
\end{proof}

The proposition says what the barrier certifies in Thorin language. The formal Thorin measure has its atoms on the positive half line, up to an unknown set of conjugate pairs which de Bruijn's theorem confines between the parabola $\mathcal P_{\sqrt{1-2t}}$ and the half line; the canopy is the parabolic band $\mathcal P_{\sqrt{1-2t}}\le\cdot<\mathcal P_{y_0}$; and the barrier is a window of heights $\eta=2xy$ in that band, cut out by two confocal left opening parabolas, through which no complex atom passes during $[0,t_0]$. It is not a statement that atoms do not leave the half line through the window: an atom that leaves the half line at a collision does so at a point of $(0,\infty)$, to the right of $\mathcal P_{y_0}$, and it is the subsequent motion into the band that the window watches. What the barrier certifies is therefore that the off axis part of $\mu_t$ to the left of the parabola $\mathcal P_{y_0}$ is not fed through a window of heights $\eta\approx2Xy$ during $[0,t_0]$, so that its fate is decided by (i$'$) on the left and by the analytic estimate (ii$'$) on the right.

\subsection{The density certificate}

In the chart $\theta=\sqrt z$, with the basepoint shift of Theorem~\ref{cp:thm:transfer}, the formal Thorin measure has the density $\pi^{-1}v_t(\sigma,\theta)$, where
\[
v_t(\sigma,\theta)=\re\frac{\xi_t'}{\xi_t}(\sigma+i\theta),
\]
the function of Section~\ref{cp:sec:minflow} and $\pi\rho_{\sigma-1/2,t}(\theta)$ of Section~\ref{sec:ch21:flow}. It is a positive Thorin density at the basepoint $\sigma$ when no zero of $\xi_t$ lies to the right of $\sigma$ (Theorem~\ref{cp:thm:transfer}); at a general $\sigma$ it is the Poisson sum of Proposition~\ref{prop:struct} for $\xi_t$, with negative kernels from the zeros to the right of $\sigma$. The barrier rectangle has a pointwise certificate in terms of this density.

\begin{proposition}[The barrier in density form]\label{barrier:prop:density}
Let $0\le t\le t_0$, $\frac12<\alpha<\sigma_1$ and $\theta_1<\theta_2$, and suppose $v_t(\sigma,\theta)>0$ at every point of $[\alpha,\sigma_1]\times[\theta_1,\theta_2]$ at which $\xi_t(\sigma+i\theta)\ne0$. Then $\xi_t$ has no zero $\beta+i\gamma$ with $\alpha<\beta\le\sigma_1$ and $\theta_1\le\gamma\le\theta_2$. In particular, if $v_t>0$ on $\mathcal B(X,y_0)$ off the zeros of $\xi_t$ for every $t\in[0,t_0]$, then $\xi_t$ has no zero in $\mathcal B(X,y_0)$ with $\beta>\frac12+\frac{y_0}2$ for any such $t$.
\end{proposition}

\begin{proof}
Near a zero $\rho$ of multiplicity $m$, $\xi_t'/\xi_t(s)=m/(s-\rho)+O(1)$, so $v_t(\beta-\delta,\gamma)=-m/\delta+O(1)<0$ for small $\delta>0$, and these points lie in the rectangle when $\alpha<\beta\le\sigma_1$ and $\theta_1\le\gamma\le\theta_2$. This is the proof of Proposition~\ref{cp:prop:thorincert} with the bottom edge $\theta=0$ replaced by $\theta_1$.
\end{proof}

Zeros on the left edge $\re s=\alpha$ are not excluded by the proposition; they are the zeros with $y=y_0$, which hypothesis (iii) in the form \eqref{barrier:eq:box} excludes and which the winding number of Section~\ref{barrier:sec:argument} also detects. The density certificate is a two dimensional condition on the whole rectangle, because positivity on one line does not exclude zeros to its right (Remark~\ref{rem:ch2:oneline}). It is stronger than the absence of zeros, since $v_t$ can be negative at points of the rectangle where no zero lies to their right within it, so it can fail where the winding test succeeds. What it offers is a margin: the minimum of $v_t$ over $\mathcal B$ is a number, comparable across times and heights, and Section~\ref{barrier:sec:check} measures it.

\section{A computed barrier at a modest height}\label{barrier:sec:check}

We ran the box \eqref{barrier:eq:box} of the barrier with $t_0=y_0=0.2$, as in \cite{Polymath}, at a height that the P\'olya route reaches in minutes, and with both certificates: the winding number of $\xi_t$ around $\partial\mathcal B$ and the minimum of $v_t$ on $\mathcal B$. The run is a check of the method and of the evaluator, not a contribution to any bound on $\Lambda$: a barrier is useful only at a height up to which RH has been verified and beyond which (ii) is known, and at $X\approx10^3$ hypothesis (ii) is not available.

\subsection{What is computed}

The check does the following.
\begin{enumerate}
\item The window. A scan of $\min|\zeta|$ over $\sigma\in\{0.6,0.8,1\}$ and five heights in each window $[\theta,\theta+\frac12]$, $495\le\theta<505$ in steps of $0.05$, at $20$ digits, chose $X/2=497.4$, where the minimum is $2.2472$. The window $[497.4,497.9]$ lies between the zeros $\gamma_{268}=496.4297$ and $\gamma_{269}=498.5808$ of $\zeta$. So $X=994.8$, and the rectangle is $\mathcal B=[0.6,1]\times[497.4,497.9]$. This imitates, on a small scale, the choice of $X$ in \cite{Polymath} that makes $|H_t|$ large on the barrier.
\item {\emergencystretch=4em The evaluator. $\xi_t(s)$, $\xi_t'(s)$ and the derivatives $\partial_t^k\xi_t(s)=8\int_0^\infty u^{2k}e^{tu^2}\allowbreak\Phi_{\rm dBN}(u)\allowbreak\cosh(2(s-\frac12)u)\,du$ for $k\le8$ are computed by the trapezoid rule of Section~\ref{sec:ch4:trapezoid} with the weight $e^{tu^2}$, on $[0,3]$ with step $h=0.002361$ ($1273$ nodes), at a working precision of $214$ digits; $h$ is chosen from the Poisson summation bound with the strip $|\im u|<0.3$, as in Section~\ref{cp:sec:march}, for $30$ correct digits relative to $|\xi_t|$, which is about $4\times10^{-165}$ on $\mathcal B$. At the four corners of $\mathcal B$ the evaluator at $t=0$ agrees with $\xi(s)=\frac12s(s-1)\pi^{-s/2}\Gamma(\frac s2)\zeta(s)$ computed by \texttt{mpmath} \cite{mpmath} to a relative error between $2.1\times10^{-51}$ and $4.0\times10^{-51}$, and at $t=0.2$ a second rule with step $0.8h$ agrees to at most $3.6\times10^{-51}$.\par}
\item The winding numbers. For each $t_i=0,0.01,\dots,0.2$ the boundary of $\mathcal B$ is walked counterclockwise with an adaptive step keeping $|\xi_{t_i}(s_{j+1})-\xi_{t_i}(s_j)|<\frac12|\xi_{t_i}(s_j)|$, as in Lemma~\ref{cp:lem:grid} but tested only at the points of the walk, and the winding number is the sum of the principal arguments of $\xi_{t_i}(s_{j+1})/\xi_{t_i}(s_j)$ divided by $2\pi$.
\item Between grid times. At every point $s_j$ of the walk at time $t_i$ we evaluate the relative Taylor sum $\sum_{k=1}^8\frac{D^k}{k!}|\partial_t^k\xi_{t_i}(s_j)|/|\xi_{t_i}(s_j)|$ with $D=t_{i+1}-t_i$, the form of (B2) of Section~\ref{sec:ft:persist} that Proposition~\ref{nw:prop:height} forces: bounds through the absolute value of the integrand cannot certify any step at an appreciable height. The rest of the Taylor series beyond order $8$ is bounded by moving the contour to $\im u=\eta$ as in Lemma~\ref{cp:lem:tilt}:
\[
\begin{aligned}\Big|\sum_{k>8}\frac{D^k}{k!}\partial_t^k\xi_t(s)\Big|\\{}&\le8e^{-2\eta\theta}\int_0^\infty e^{t(u^2-\eta^2)}\\{}&\qquad{}\cdot|\Phi_{\rm dBN}(u+i\eta)|\cosh(2xu)\frac{r^9e^r}{9!}\,du,\qquad r=D|u+i\eta|^2,\end{aligned}
\]
for $s=\frac12+x+i\theta$, by $|(u+i\eta)^{2k}|=|u+i\eta|^{2k}$ and $\sum_{k>8}r^k/k!\le r^9e^r/9!$; the integral is computed by Gauss Legendre quadrature at $x=\frac12$, $\theta=497.4$ and $\eta=0.38$, which bounds it on $\mathcal B$ because the integrand increases with $x$ and $e^{-2\eta\theta}$ decreases with $\theta$.
\item The density. The minimum of $v_t$ over a grid of $9\times11$ points of $\mathcal B$ (steps $0.05$ in $\sigma$ and in $\theta$) at each $t_i$.
\item A control. At $t=0$ the walk is repeated for $\xi(s)\,(w^2-w_0^2)(w^2-\bar w_0^2)$, $w=s-\frac12$, $w_0=\frac14+i(X/2+\frac14)$, which has exactly one zero, $0.75+497.65\,i$, in $\mathcal B$.
\end{enumerate}

\subsection{The results}

\begin{table}[ht]
\centering\small
\begin{tabular}{@{}lccccc@{}}
\toprule
$t$ & winding & points & $\min|\xi_t|$ on the walk & Taylor ratio, $D=0.01$ & $\min_{\mathcal B}v_t$\\
\midrule
$0$ & $-2.6\times10^{-217}$ & $35$ & $4.087\times10^{-165}$ & $0.00773$ & $0.25392$\\
$0.05$ & $-1.3\times10^{-216}$ & $35$ & $4.123\times10^{-165}$ & $0.00781$ & $0.25815$\\
$0.10$ & $4.6\times10^{-217}$ & $35$ & $4.155\times10^{-165}$ & $0.00788$ & $0.26252$\\
$0.15$ & $-1.5\times10^{-216}$ & $35$ & $4.184\times10^{-165}$ & $0.00795$ & $0.26705$\\
$0.19$ & $5.8\times10^{-218}$ & $35$ & $4.204\times10^{-165}$ & $0.00800$ & $0.27064$\\
$0.20$ & $-1.9\times10^{-216}$ & $35$ & $4.208\times10^{-165}$ & & $0.27145$\\
\bottomrule
\end{tabular}
\caption{The barrier $\mathcal B=[0.6,1]\times[497.4,497.9]$, $X=994.8$, $y_0=0.2$, at six of the $21$ grid times (computed, floating point at $214$ digits; the rows $t=0.05$ to $0.19$ are not recomputed in ball arithmetic). Taylor ratio: the largest over the walk of $\sum_{k\le8}\frac{D^k}{k!}|\partial_t^k\xi_t|/|\xi_t|$.}
\label{barrier:tab:check}
\end{table}

\begin{proposition}[Computed, and certified in ball arithmetic for all $t\in{[0,0.2]}$]\label{barrier:prop:check}
With $t_0=y_0=0.2$ and the rectangle $\mathcal B=[0.6,1]\times[497.4,497.9]$:
\begin{enumerate}
\item[(a)] at each of the $21$ times $t_i=0,0.01,\dots,0.2$ the computed winding number of $\xi_{t_i}$ around $\partial\mathcal B$ is $0$, the sums of arguments being below $3\times10^{-216}$ in absolute value, with $35$ points per walk and the ratio $|\Delta\xi|/|\xi|$ at most $0.169$ along the walks; the control function returns winding number $1.0$ in $69$ points;
\item[(b)] at every point of every walk the relative Taylor sum of order $8$ over the step $D=0.01$ is at most $0.00801$, and the tilted tail beyond order $8$ is at most $7.4\times10^{-186}$, that is $1.8\times10^{-21}$ relative to the smallest $|\xi_t|$ met; one Taylor step from $t=0$ over the whole interval, $D=0.2$, gives the relative sum $0.171$ and a relative tail below $9.9\times10^{-10}$;
\item[(c)] the minimum of $v_t$ over the $9\times11$ grid of $\mathcal B$ increases with $t$, from $0.25392$ at $t=0$ to $0.27145$ at $t=0.2$, and it is attained on the left edge $\sigma=0.6$, at $\theta=497.5$ for $t\le0.17$ and at $\theta=497.45$ for $t\ge0.18$.
\end{enumerate}
The floating point values above agree with a ball arithmetic computation, which certifies for every $t\in[0,0.2]$ (not only at the grid times): the winding number of $\xi_0$ around $\partial\mathcal B$ is $0$ (and the control returns $1$); one Taylor step of order $8$ from $t=0$ over $D=0.2$ gives $|\xi_t-\xi_0|\le0.198|\xi_0|$ on all of $\partial\mathcal B$ for all $t\in[0,0.2]$, hence $\xi_t$ has no zero on $\partial\mathcal B$, the winding number stays $0$, and there is no zero of $\xi_t$ in $\mathcal B$ for $t\in[0,0.2]$; and $v_t\ge0.1587$ on all of $\mathcal B$ for all $t\in[0,0.2]$. Still floating point or taken from elsewhere: the table rows at $t=0.05$ to $0.19$, the untilted bound below, the monotone increase of $\min v_t$ in $t$ in part (c) (the certified bound is $0.1587$, below the grid values because the cells are discs), and the Polymath15 barrier hypotheses (i) and (ii) of Theorem~\ref{barrier:thm:pm} and their use at $X=994.8$.
\end{proposition}

Read with Section~\ref{barrier:sec:argument}, part (a) says that at the grid times no zero of $\xi_t$ lies in $\mathcal B$, and part (b) says that between grid times $|\xi_t-\xi_{t_i}|<\frac12|\xi_{t_i}|$ at the points of the walks, which by Rouch\'e's theorem would carry the winding number $0$ across each interval if it held on the whole of $\partial\mathcal B$. Part (c) gives the same conclusion pointwise through Proposition~\ref{barrier:prop:density}, again only on a grid. The untilted bound for the same tail, the absolute value bound of Proposition~\ref{nw:prop:height}, is $1.4\times10^{-23}$ for $D=0.2$, some $10^{141}$ times $|\xi_t|$ on $\mathcal B$, so the tilt is essential even at this height.

Three gaps separated the floating point proposition from a certified barrier at this height, and the ball arithmetic computation closes each by the methods of Section~\ref{cp:sec:cert}: the evaluator runs in ball arithmetic with the truncation and discretisation errors in the radius (Theorem~\ref{cp:thm:rig}); the winding at $t=0$ is verified on the segments between walk points through a derivative bound (Lemmas~\ref{cp:lem:deriv} and~\ref{cp:lem:grid}); and the Taylor criterion is verified on all of $\partial\mathcal B$ by the same derivative bound applied to $\partial_t^k\xi_t$, with the tail integral enclosed by interval subdivision, as the constants $\bar K_\eta$ are in Section~\ref{cp:sec:march}. The certified bounds are weaker than the point values (the one step ratio is $0.198$ against $0.171$, the supremum over discs against the maximum over the walk), and the threshold $1$ of Rouch\'e's theorem is met with a factor above $5$ to spare.

Two features of the numbers are worth recording. First, the barrier is cheap in $t$: at this height a single Taylor step covers $[0,0.2]$ with a factor of almost $3$ to spare against the threshold $\frac12$ used in the walks (Rouch\'e's theorem needs only $1$), whereas \cite{Polymath} needed many time values at $X\approx6\times10^{10}$; heuristically, the relative size of $\partial_t\xi_t=\frac14\partial_w^2\xi_t$ is of order $(\log\theta)^2$ near zeros and of order the inverse squared distance to the nearest zero, so the number of time steps grows with the height and with the closeness of the zeros to the window. Second, $\min v_t$ on $\mathcal B$ moves by only $7\%$ between $t=0$ and $t=0.2$, in line with Proposition~\ref{cp:prop:see}(iii): the density certificate does not see the difference between the time at which the clock is known to be a GGC and the time at which it is not known to be one.

\section{What the Thorin form adds, and why it cannot reach \texorpdfstring{$t=0$}{t=0}}\label{barrier:sec:limits}

\subsection{What it adds}

Three things. First, a single statement that joins the method to the rest of the book: the output of the barrier method is a point $(t_0,\frac12+\frac{y_0}2)$ of the GGC region $\mathcal G$ (Proposition~\ref{barrier:prop:ggc}), and every later improvement is a point further down the phase diagram of Chapter~\ref{chap:phase}. Second, a geometry: the canopy, the barrier and the strip theorem become a parabolic band, a window between confocal parabolas, and a family of right opening parabolas collapsing onto the half line (Lemma~\ref{barrier:lem:quad}, Proposition~\ref{barrier:prop:window}), the same picture as the shrinking parabolas of Corollary~\ref{ki:cor:parab}. Third, a quantitative certificate: the density $v_t$ on the barrier rectangle (Proposition~\ref{barrier:prop:density}) gives a margin that can be compared across times and heights, which the winding number, an integer, does not.

\subsection{What it does not add}

It adds no analytic input. Hypothesis (ii) is the heart of the method, and the Thorin form restates it as (ii$'$) without proving it; Polymath15 prove it through an effective Riemann Siegel approximation of $H_t$, and the Thorin side has no substitute. Lemma~\ref{cp:lem:shift} writes $\xi_t$ as a Gaussian average of $\zeta$ shifted to the right by about $\frac t4\log\frac\theta{2\pi}$, which is one way to see why a positive time helps: the average samples $\zeta$ at a real part that grows with the height, where the Euler product dominates. The density certificate is also weaker than the winding test in the sense of Section~\ref{barrier:sec:thorin}: it can fail where no zero lies, and its margin at basepoints near $\frac12$ tends to zero (Remark~\ref{rem:dbeta:margin}), so it can only be used with $y_0$ bounded away from $0$.

\subsection{Why the method cannot reach \texorpdfstring{$t=0$}{t=0}}

The bound of Theorem~\ref{barrier:thm:pm} is $t_0+\frac12y_0^2$, and RH is $\Lambda\le0$, so a proof of RH by the method needs $t_0\to0$ and $y_0\to0$. Each limit fails for a separate reason.

As $t_0\to0$, hypothesis (ii) becomes a zero free region for $\zeta$ itself: no zero with $\re s\ge\frac12+\frac{y_0}2$ above height $X/2$. For $y_0<1$ no such region is known; the regions of Vinogradov and Korobov narrow towards $\re s=1$, and the half plane $\re s>\frac78$ announced in October 2026 and not refereed \cite{OAI26a,OAI26b,Liu26a} would give it only for $y_0\ge\frac34$, that is the point $(0,\frac78)$ of Corollary~\ref{cp:cor:known}(d) and the bound $\Lambda\le\frac9{32}$, worse than $0.2$. The effective approximation that proves (ii) at positive time depends on the time: heuristically the flow at time $t$ makes the zeros near height $x$ real once $t\log x$ exceeds a constant, and the theorem of \cite{Polymath} cited in the proof of Corollary~\ref{ft:cor:collide} says that for an absolute constant $C$ all zeros of $H_t$ with $|x|\ge e^{C/t}$ are real and simple (we read only its statement as the book cites it, not its proof). The required height $X$ therefore grows like $e^{c/t_0}$, and the table of \cite{Polymath} fits $\log X\approx4.24/\Lambda+8.28$ (Proposition~\ref{cp:prop:cost}), which needs verification of RH to about $10^{40}$ for $\Lambda\le0.05$. In the limit $t_0=0$ the barrier height is infinite, and the method has reduced RH to RH.

As $y_0\to0$, the barrier rectangle $\mathcal B$ widens to the left edge $\re s=\frac12$, on which $\xi_t$ has its real zeros for $t\ge\Lambda$, and the band of Lemma~\ref{barrier:lem:quad} closes onto the half line, where the atoms of $\mu_t$ live. The winding test on a rectangle touching the critical line needs a window free of zeros on the line at every time in $[0,t_0]$, which moving zeros make hard but not impossible (the box barrier of Proposition~\ref{ft:prop:below}, not carried out; Remark~\ref{mw:rem:flowcert}), and the density margin vanishes (Remark~\ref{rem:dbeta:margin}). More fundamentally, a statement at $y_0=0$ and $t_0>0$ would only give $\Lambda\le t_0$.

Finally, the Thorin form cannot make up the difference, for the reason given in Section~\ref{ov1:sec:path}: a perturbation of $\xi$ with the same zeros below any height passes the finite tests of the heat trace listed in Theorem~\ref{thm:ch13:perturb}, so a finite barrier computation together with an analytic tail at $t_0>0$ can certify points of $\mathcal G$ at positive time but never the corner $(0,\frac12)$. By Rodgers and Tao, $\Lambda\ge0$ \cite{RodgersTao}, so the corner is on the boundary of the region the method can reach, and the step to it is a closed condition at $t=0$ that must come from an analytic argument which uses the Euler product (Proposition~\ref{ki:prop:obstruction} rules out arguments that survive a Gaussian tilt).

\section{Open problems}\label{barrier:sec:open}

\begin{problem}\label{barrier:prob:certify}
Proposition~\ref{barrier:prop:check} is certified in ball arithmetic at $X=994.8$ (Section~\ref{barrier:sec:check}). Determine, as a function of the height $X$, the number of time steps that a certified barrier of width one at $t_0=y_0=0.2$ needs. Is it of order $(\log X)^2$, as the size of $\partial_t\xi_t/\xi_t$ near the zeros suggests?
\end{problem}

\begin{problem}\label{barrier:prob:margin}
Measure $\min_{\mathcal B}v_t$ on barrier rectangles at heights up to $10^{10}$ with a Riemann Siegel type evaluator (step 1 of Section~\ref{cp:sec:programme}), at the windows chosen by the Euler product heuristic of \cite{Polymath}. Does the density certificate succeed wherever the winding test does, for $y_0=0.2$, and how does its margin depend on the distance from the window to the nearest zeros?
\end{problem}

\begin{problem}\label{barrier:prob:window}
In the chart of squares, the barrier is a window between $\mathcal Q_X$ and $\mathcal Q_{X+1}$. Is there a version of Theorem~\ref{barrier:thm:pm} in which the window is replaced by a condition on the Thorin measure alone, for instance a bound on the mass of $\mu_t$ near $\mathcal Q_X$ along the flow, that is equivalent to (iii) and can be checked from moments of the P\'olya kernel tilted by $e^{tu^2}$?
\end{problem}

\begin{problem}\label{barrier:prob:tail}
Find an analytic statement of type (ii) at $t_0=0$ and some $y_0<\frac34$ (the announced half plane would give $y_0=\frac34$ conditionally) that uses the Euler product: a zero free region for $\zeta$ to the right of $\re s=\frac12+\frac{y_0}2$ above a computable height. With the barrier and the verification of RH it would give $\Lambda\le\frac12y_0^2$; it is a quasi Riemann hypothesis, and the problem is the hard core identified in Section~\ref{ov1:sec:path}.
\end{problem}

\subsection*{Does the barrier method give a path to RH?}

No. It gives a monotone sequence of theorems $\Lambda\le c$ with $c>0$, each of them a point of the GGC region at positive time, and its cost grows like $e^{c'/\Lambda}$ in the height to which RH must be verified. The Thorin form states each step as a GGC theorem and adds a parabolic geometry and a density margin, but it supplies no analytic input below the basepoint one, and the limit $t_0\to0$ turns the analytic hypothesis of the method into a zero free region for $\zeta$, which is the open step of the path to RH and not a way around it.
